\documentclass[12pt]{article}
\usepackage[utf8]{inputenc}
\usepackage[T1]{fontenc}
\usepackage{hyperref}
\title{Slovenian Sign Language Language Resources for K-12 Curriculum Development}
\author{Compiled by Claude Sonnet 5.5 (claude-sonnet-5-5) for Andr\'as Kornai}
\date{}
\begin{document}
\maketitle

\section{Introduction}
Slovenian Sign Language (slovenski znakovni jezik, SZJ; ISO 639-3: ysl) is the sign language of the Slovenian deaf community. The search results say it was written into the Slovenian Constitution in 2021. They give the size of its lexicon as about 25,000 signs, a figure taken from a campaign statement and not independently counted. Several sources say the language is poorly documented linguistically and that school teaching materials have been inadequate.

\section{Primary Educational Resources}
\begin{itemize}
\item Ministry of Education curriculum for the elective subject ``Slovenski znakovni jezik'' in primary school: 204 hours across grades 7--9, in adapted programmes. \url{https://www.gov.si/assets/ministrstva/MIZS/Dokumenti/Osnovna-sola/Ucni-nacrti/izbirni/3-letni-lahko-krajsi/Slovenski_znakovni_jezik_izbirni.pdf}
\item Ministry news item (2 June 2023) on better conditions for SZJ users in education. \url{https://www.gov.si/novice/2023-06-02-za-boljse-moznosti-uporabnikov-slovenska-znakovnega-jezika-v-vzgoji-in-izobrazevanju/}
\item Ministry event page: a public consultation on SZJ users' needs in education, listed for 20 October 2026, which is upcoming. \url{https://www.gov.si/dogodki/2026-10-20-javno-posvetovanje-o-potrebah-uporabnikov-znakovnega-jezika-pri-vzgojno-izobrazevalnem-delu/}
\end{itemize}

\section{Linguistic and Reference Materials}
\begin{itemize}
\item LREC sign-language workshop paper (Hamburg archive) on Slovenian Sign Language resources. The title and exact content were not visible in the search result. \url{https://www.sign-lang.uni-hamburg.de/lrec/pub/16019.pdf}
\item Article in Slovenski Jezik / Slovene Linguistic Studies (srl.si) reporting on lexicography for SZJ. \url{https://srl.si/ojs/srl/article/view/4089/3560}
\item Journal of Language and Education, December 2023 (DOAJ entry), on video reference grammar materials for SZJ. \url{https://doaj.org/article/1be6afd6904d47e7bfa8168834f5d117}
\item Journal article on education (University of Maribor journals). Its topic could not be confirmed from the snippet. \url{https://journals.um.si/index.php/education/article/view/2986/2285}
\item INSIGN national report for Slovenia. \url{https://www.uni-siegen.de/zew/insign/outcomes/communication_challenges/files/national_report_-_insign_-_slovenia_-_final.pdf}
\item Proceedings paper from ISJT 2014. Its topic could not be confirmed from the snippet. \url{https://nl.ijs.si/isjt14/proceedings/isjt2014_29.pdf}
\end{itemize}

\section{Religious and Community Resources}
\begin{itemize}
\item The Deaf and Hard of Hearing Clubs Association is named in news reports as running a crowdfunding campaign for new signs with Mastercard Slovenia. See the Slovenia Times report. \url{https://sloveniatimes.com/41453/slovenian-sign-language-getting-a-boost}
\item STA report on an online platform launched to develop Slovenian sign language. \url{https://english.sta.si/3362027/online-platform-launched-to-develop-slovenian-sign-language}
\item RTV Slovenija feature: sign language is a language, not a tool for communicating with a deaf person. \url{https://www.rtvslo.si/dostopno/clanki/znakovni-jezik-je-jezik-in-ne-orodje-za-komunikacijo-z-gluho-osebo/507749}
\end{itemize}

\section{Digital and Audio Resources}
\begin{itemize}
\item An online and mobile SZJ dictionary was built by the University of Maribor and NUIMO with the Slovenian Association of the Deaf and Hard of Hearing. The search results described it but gave no direct URL, so none is listed here. A conference poster from SIRikt 2013 relates to this area. \url{https://skupnost.sio.si/sio_arhiv/sirikt/www.sirikt.si/slo/prispevki/plakati/947.html}
\end{itemize}

\section{Curriculum Development Context}
SZJ now has constitutional status and a ministry action plan (September 2021) for placing it in preschool and primary and secondary education. Reports still describe a lack of linguistically grounded teaching materials. The existing dictionary, the video grammar work and the elective-subject curriculum are the main foundations. Full K-12 instruction through SZJ remains to be built.
\end{document}